\begin{proposition}\label{sga3:XXI:4.1}
\footnote{The following numbers \(4.1\) to \(4.4\) are used in the
proof of Theorem~\ref{sga3:XXI:5.1}. There are now simpler proofs of
that theorem; see \emph{Groupes et algèbres de Lie}, Chapter~V,
\S3, Theorem~1. Editors' note: The reference was made explicit and
this note, which occurred at \(4.2\) in the original, was moved here.}
Let \(n\) be the order of \(s_\alpha s_\beta\). Put \(u_0=0\), and,
for \(p\in\mathbb N\), define
\[
\left\{
\begin{aligned}
u_{2p+1}&=u_{2p}+(s_\alpha s_\beta)^p(\alpha),\\
u_{2p+2}&=u_{2p+1}+(s_\alpha s_\beta)^p s_\alpha(\beta).
\end{aligned}
\right.
\]
Then:
\footnote{Editors' note: Let \(\rho\) be the half-sum of the positive
roots, cf. Proposition~\ref{sga3:XXI:3.5.1}. If, as below, one puts
\(w_0=1\), \(w_1=s_\alpha\), \(w_2=s_\beta s_\alpha\),
\(w_3=s_\alpha s_\beta s_\alpha\), and so on, then
\(u_k=\rho-w_k^{-1}(\rho)\). This proves \SGARefLink{sga3:XXI:4.1:item:i}{(i)}, since \(w_{2n}=1\).}
\begin{enumerate}
\item[\SGAInlineRefTarget{sga3:XXI:4.1:item:i}(i)] \(u_{k+2n}=u_k\), for every \(k\in\mathbb N\).
\item[\SGAInlineRefTarget{sga3:XXI:4.1:item:ii}(ii)] \(u_0=0\), \(u_1=\alpha\), \(u_{2n-1}=\beta\),
and \(u_{2n}=0\).
\item[(iii)] If \(1<k<2n-1\), then
\[
 u_k=a_k\alpha+b_k\beta,
 \qquad a_k,b_k\in\mathbb N^\ast.
\]
\end{enumerate}
\end{proposition}

\begin{proof}
Assertion \SGARefLink{sga3:XXI:4.1:item:i}{(i)} follows from \((s_\alpha s_\beta)^n=1\) and
\(u_{2n}=0\).

We prove \SGARefLink{sga3:XXI:4.1:item:ii}{(ii)} and \SGARefLink{sga3:XXI:3.6:item:iii}{(iii)}. Direct calculation gives, in the four cases,
the following sequences of values:
\footnote{Editors' note: The original was corrected here to agree
with the convention \(\ell(\alpha)\le\ell(\beta)\) of \(4.0\).}
\[
\begin{array}{c|l}
A_1\times A_1&0,\alpha,\beta+\alpha,\beta,0,\\[2pt]
A_2&0,\alpha,\beta+2\alpha,2\beta+2\alpha,
       2\beta+\alpha,\beta,0,\\[2pt]
B_2&0,\alpha,\beta+3\alpha,2\beta+4\alpha,
       3\beta+4\alpha,3\beta+3\alpha,2\beta+\alpha,\beta,0,\\[2pt]
G_2&0,\alpha,\beta+4\alpha,2\beta+6\alpha,
       4\beta+9\alpha,5\beta+10\alpha,6\beta+10\alpha,\\
   &\quad 6\beta+9\alpha,5\beta+6\alpha,4\beta+4\alpha,
       2\beta+\alpha,\beta,0.
\end{array}
\]
\end{proof}

\begin{lemma}\label{sga3:XXI:4.2}
Put
\[
 w_{2p}=(s_\beta s_\alpha)^p,
 \qquad
 w_{2p+1}=s_\alpha(s_\beta s_\alpha)^p,
\]
so that \(w_0,\ldots,w_{2n-1}\) are the distinct elements of \(W\).
Let \(u_0,\ldots,u_{2n-1}\) be the elements of \(V\) defined in
Proposition~\ref{sga3:XXI:4.1}. For every \(x\in V^\ast\), one has
\[
 x-w_k(x)=n_k\alpha^\ast+m_k\beta^\ast,
\]
where \(n_k,m_k\in\mathbb Q\) and
\[
 n_k+m_k=(x,u_k).
\]
\footnote{Editors' note: In view of the preceding editor's note, this
follows immediately from Propositions~\ref{sga3:XXI:3.5.1} and
\ref{sga3:XXI:1.1.9}: one has
\[
 n_k+m_k=(x-w_k(x),\rho)
 =(x,\rho-w_k^{-1}(\rho))=(x,u_k).
\]}
\end{lemma}

\begin{proof}
The proof is by induction on \(k\). If \(k=0\), the formula is
trivial. We carry out, for example, the passage from \(w_{2p}\) to
\(w_{2p+1}\). Since \(w_{2p+1}=s_\alpha w_{2p}\),
\[
\begin{aligned}
x-w_{2p+1}(x)
 &=x-w_{2p}(x)+w_{2p}(x)-s_\alpha w_{2p}(x)\\
 &=n_{2p}\alpha^\ast+m_{2p}\beta^\ast
   +(w_{2p}(x),\alpha)\alpha^\ast.
\end{aligned}
\]
Consequently
\[
\begin{aligned}
n_{2p+1}+m_{2p+1}
 &=n_{2p}+m_{2p}
   +\bigl((s_\beta s_\alpha)^p(x),\alpha\bigr)\\
 &=(x,u_{2p})+\bigl(x,(s_\alpha s_\beta)^p(\alpha)\bigr)\\
 &=(x,u_{2p+1}).
\end{aligned}
\]
\end{proof}
